\chapter{La movilidad en Bogotá: datos y proyectos}\label{cap:contexto}

\section{Demanda de viajes}

La Encuesta de Movilidad 2023 estima 14,64 millones de viajes diarios en Bogotá y veinte municipios
vecinos, de los cuales 12,14 millones se originan en Bogotá y 2,5 millones en los municipios
\citep{sdm2024}. El transporte público concentra el 35\,\% de los viajes, la caminata el 28\,\%, el
automóvil el 14\,\%, la bicicleta y la motocicleta el 7\,\% cada una y el taxi el 4\,\%; la duración
promedio de un viaje es de 48 minutos. Entre los modos motorizados que compiten por la capacidad de los
corredores (transporte público, automóvil, motocicleta y taxi, que suman alrededor del 60\,\% de los
viajes), el transporte público representa entonces cerca de 35/60 $\approx$ 58\,\% de los viajes.

La estructura espacial es marcadamente radial: la generación de viajes se concentra en las localidades
periféricas y la atracción en el centro ampliado y el corredor oriental, lo que produce flujos de hora
pico en un solo sentido sobre los corredores troncales. Según las proyecciones posteriores a la
pandemia, Bogotá tiene 7.937.898 habitantes en 2025 \citep{dane2025}; Suba (cerca de 1,2 millones) y
Kennedy (cerca de 1,1 millones) son las localidades más pobladas, y la población total empezaría a
decrecer desde 2026, con crecimiento solo en localidades del borde como Bosa, Usme y Ciudad Bolívar.
Los municipios de la Sabana, en cambio, crecen con rapidez: Soacha tiene cerca de 870 mil habitantes y
Mosquera y Chía superan los 170 mil cada uno \citep{dane2025}.

\section{Oferta y desempeño}

TransMilenio movió cerca de 2 millones de pasajeros troncales por día hábil en 2025, todavía por
debajo de los 2,6 millones de 2018--2019; la velocidad comercial de las troncales fue de 25,2 km/h en
2024 \citep{silla2026}. En la red vial, la velocidad media en hora pico fue de 18,9 km/h en 2025, lo que
ubicó a Bogotá en el séptimo lugar mundial de congestión \citep{tomtom2026}.

\section{Proyectos en curso y en decisión}

La \cref{tab:proyectos} resume los proyectos que considera el modelo. Se distinguen los
\emph{comprometidos}, contratados o en construcción, que forman parte de todos los escenarios, de los
\emph{en decisión}, que definen los diseños urbanos. Los costos de los proyectos en decisión son, salvo
la extensión de la Línea 1, estimaciones del autor a partir de cifras públicas y se tratan con
sobrecostos aleatorios (\cref{sec:costo}).

\begin{table}[ht]
\centering\small
\caption{Proyectos de la red modelada (CAPEX en billones de pesos de 2025).}\label{tab:proyectos}
\resizebox{\textwidth}{!}{\begin{tabular}{llllrrr}
\toprule
Id & Proyecto & Estado & Apertura & CAPEX & Capacidad & Velocidad \\
 & & & & billones & pas/h-sentido & km/h \\
\midrule
L1 & Metro Línea 1 (Bosa – Calle 72) & comprometido & 2028 & — & 72.000 & 43 \\
T68 & Troncal Av. 68 (TransMilenio) & comprometido & 2027 & — & 25.000 & 25 \\
RTO & Regiotram de Occidente & comprometido & 2030 & — & 10.000 & 30 \\
L1X & Extensión Línea 1 hasta la Calle 100 & en decisión & 2031 & 3,4 & 72.000 & 43 \\
L2 & Metro Línea 2 (Calle 72 – Suba) & en decisión & 2033 & 19,6 & 60.000 & 40 \\
L3 & Metro Línea 3 (Soacha – Calle 100) & en decisión & 2033 & 31,0 & 60.000 & 40 \\
RTN & Regiotram del Norte & en decisión & 2034 & 4,3 & 14.000 & 35 \\
\bottomrule
\end{tabular}
}
\fuente{Empresa Metro de Bogotá \citep{emb2026,emb2026l2,emb2026l3}; Gobernación de Cundinamarca
\citep{regiotram2026}; IDU; cálculos del autor. El costo de la Línea 2 corresponde a 16,25 billones de
pesos de noviembre de 2022 \citep{emb2023} indexados por \Indexacion; el de la Línea 3 aplica el costo unitario de la Línea 2 a sus 24,5 km; el del
Regiotram del Norte aplica el costo unitario del Regiotram de Occidente a unos 48 km.}
\end{table}

\section{Ordenamiento territorial}

El Plan de Ordenamiento Territorial adoptado con el Decreto 555 de 2021 propone una ciudad de
proximidad (la «ciudad de 30 minutos»), organizada en unidades de planeamiento local y con
centralidades que acerquen empleo y servicios a la vivienda \citep{decreto555}. En este trabajo esa
orientación se representa con una política de intensidad $\theta\in[0,1]$ que reduce la fracción de
viajes motorizados y relocaliza parte de la atracción del centro hacia subcentros asociados a la red
de transporte masivo (\cref{sec:suelo}).
